\begin{register}{H}{I2C\_SCL\_LOW\_PERIOD\_REG}{\hexprefix{}0000}\label{regdesc:I2CSCLLOWPERIODREG}
\regfield{(reserved)}{18}{14}{000000000000000000}%
\regfield{I2C\_SCL\_LOW\_PERIOD}{14}{0}{{\hexprefix{}00}}%
\reglabel{Reset}\regnewline%
\begin{regdesc}\begin{reglist}
\label{fielddesc:I2CSCLLOWPERIOD}\item [I2C\_SCL\_LOW\_PERIOD] 用于配置 SCL 低电平的保持时间，以 I2C 模块时钟周期数为单位。 (读/写)
\end{reglist}\end{regdesc}
\end{register}


\begin{register}{H}{I2C\_SDA\_HOLD\_REG}{\hexprefix{}0030}\label{regdesc:I2CSDAHOLDREG}
\regfield{(reserved)}{22}{10}{0000000000000000000000}%
\regfield{I2C\_SDA\_HOLD\_TIME}{10}{0}{{\hexprefix{}0}}%
\reglabel{Reset}\regnewline%
\begin{regdesc}\begin{reglist}
\label{fielddesc:I2CSDAHOLDTIME}\item [I2C\_SDA\_HOLD\_TIME] 用于配置 SCL 下降沿后的数据保持时间，以 I2C 模块时钟周期数为单位。 (读/写)
\end{reglist}\end{regdesc}
\end{register}


\begin{register}{H}{I2C\_SDA\_SAMPLE\_REG}{\hexprefix{}0034}\label{regdesc:I2CSDASAMPLEREG}
\regfield{(reserved)}{22}{10}{0000000000000000000000}%
\regfield{I2C\_SDA\_SAMPLE\_TIME}{10}{0}{{\hexprefix{}0}}%
\reglabel{Reset}\regnewline%
\begin{regdesc}\begin{reglist}
\label{fielddesc:I2CSDASAMPLETIME}\item [I2C\_SDA\_SAMPLE\_TIME] 用于配置采样 SDA 的时间，以 I2C 模块时钟周期数为单位。 (读/写)
\end{reglist}\end{regdesc}
\end{register}


\begin{register}{H}{I2C\_SCL\_HIGH\_PERIOD\_REG}{\hexprefix{}0038}\label{regdesc:I2CSCLHIGHPERIODREG}
\regfield{(reserved)}{4}{28}{0000}%
\regfield{I2C\_SCL\_WAIT\_HIGH\_PERIOD}{14}{14}{{\hexprefix{}00}}%
\regfield{I2C\_SCL\_HIGH\_PERIOD}{14}{0}{{\hexprefix{}00}}%
\reglabel{Reset}\regnewline%
\begin{regdesc}\begin{reglist}
\label{fielddesc:I2CSCLHIGHPERIOD}\item [I2C\_SCL\_HIGH\_PERIOD] 用于配置 SCL 在主机模式下保持高电平的时间，以 I2C 模块时钟周期数为单位。 (读/写)
\label{fielddesc:I2CSCLWAITHIGHPERIOD}\item [I2C\_SCL\_WAIT\_HIGH\_PERIOD] 用于配置 SCL\_FSM 等待 SCL 在主机模式下翻转至高电平的时间，以 I2C 模块时钟周期数为单位。 (读/写)
\end{reglist}\end{regdesc}
\end{register}


\begin{register}{H}{I2C\_SCL\_START\_HOLD\_REG}{\hexprefix{}0040}\label{regdesc:I2CSCLSTARTHOLDREG}
\regfield{(reserved)}{22}{10}{0000000000000000000000}%
\regfield{I2C\_SCL\_START\_HOLD\_TIME}{10}{0}{{8}}%
\reglabel{Reset}\regnewline%
\begin{regdesc}\begin{reglist}
\label{fielddesc:I2CSCLSTARTHOLDTIME}\item [I2C\_SCL\_START\_HOLD\_TIME] 配置 START 命令产生时 SDA 下降沿 和 SCL 下降沿的间隔时间，以 I2C 模块时钟周期数为单位。 (读/写)
\end{reglist}\end{regdesc}
\end{register}


\begin{register}{H}{I2C\_SCL\_RSTART\_SETUP\_REG}{\hexprefix{}0044}\label{regdesc:I2CSCLRSTARTSETUPREG}
\regfield{(reserved)}{22}{10}{0000000000000000000000}%
\regfield{I2C\_SCL\_RSTART\_SETUP\_TIME}{10}{0}{{8}}%
\reglabel{Reset}\regnewline%
\begin{regdesc}\begin{reglist}
\label{fielddesc:I2CSCLRSTARTSETUPTIME}\item [I2C\_SCL\_RSTART\_SETUP\_TIME] 配置 RESTART 命令产生时 SCL 上升沿和 SDA 下降沿的间隔时间，以 I2C 模块的时钟周期数为单位。 (读/写)
\end{reglist}\end{regdesc}
\end{register}


\begin{register}{H}{I2C\_SCL\_STOP\_HOLD\_REG}{\hexprefix{}0048}\label{regdesc:I2CSCLSTOPHOLDREG}
\regfield{(reserved)}{18}{14}{000000000000000000}%
\regfield{I2C\_SCL\_STOP\_HOLD\_TIME}{14}{0}{{\hexprefix{}00}}%
\reglabel{Reset}\regnewline%
\begin{regdesc}\begin{reglist}
\label{fielddesc:I2CSCLSTOPHOLDTIME}\item [I2C\_SCL\_STOP\_HOLD\_TIME] 配置 STOP 命令后的延迟，以 I2C 模块时钟周期数为单位。 (读/写)
\end{reglist}\end{regdesc}
\end{register}


\begin{register}{H}{I2C\_SCL\_STOP\_SETUP\_REG}{\hexprefix{}004C}\label{regdesc:I2CSCLSTOPSETUPREG}
\regfield{(reserved)}{22}{10}{0000000000000000000000}%
\regfield{I2C\_SCL\_STOP\_SETUP\_TIME}{10}{0}{{\hexprefix{}0}}%
\reglabel{Reset}\regnewline%
\begin{regdesc}\begin{reglist}
\label{fielddesc:I2CSCLSTOPSETUPTIME}\item [I2C\_SCL\_STOP\_SETUP\_TIME] 配置 SCL 上升沿和 SDA 上升沿的间隔时间，以 I2C 模块时钟周期数为单位。 (读/写)
\end{reglist}\end{regdesc}
\end{register}


\begin{register}{H}{I2C\_SCL\_ST\_TIME\_OUT\_REG}{\hexprefix{}0098}\label{regdesc:I2CSCLSTTIMEOUTREG}
\regfield{(reserved)}{8}{24}{00000000}%
\regfield{I2C\_SCL\_ST\_TO}{24}{0}{{\hexprefix{}0100}}%
\reglabel{Reset}\regnewline%
\begin{regdesc}\begin{reglist}
\label{fielddesc:I2CSCLSTTO}\item [I2C\_SCL\_ST\_TO] SCL\_FSM 状态不变的阈值。 (读/写)
\end{reglist}\end{regdesc}
\end{register}


\begin{register}{H}{I2C\_SCL\_MAIN\_ST\_TIME\_OUT\_REG}{\hexprefix{}009C}\label{regdesc:I2CSCLMAINSTTIMEOUTREG}
\regfield{(reserved)}{8}{24}{00000000}%
\regfield{I2C\_SCL\_MAIN\_ST\_TO}{24}{0}{{\hexprefix{}0100}}%
\reglabel{Reset}\regnewline%
\begin{regdesc}\begin{reglist}
\label{fielddesc:I2CSCLMAINSTTO}\item [I2C\_SCL\_MAIN\_ST\_TO] SCL\_MAIN\_FSM 状态不变的阈值。 (读/写)
\end{reglist}\end{regdesc}
\end{register}


\begin{register}{H}{I2C\_CTR\_REG}{\hexprefix{}0004}\label{regdesc:I2CCTRREG}
\regfield{(reserved)}{20}{12}{00000000000000000000}%
\regfield{I2C\_REF\_ALWAYS\_ON}{1}{11}{{1}}%
\regfield{I2C\_FSM\_RST}{1}{10}{{0}}%
\regfield{I2C\_ARBITRATION\_EN}{1}{9}{{1}}%
\regfield{I2C\_CLK\_EN}{1}{8}{{0}}%
\regfield{I2C\_RX\_LSB\_FIRST}{1}{7}{{0}}%
\regfield{I2C\_TX\_LSB\_FIRST}{1}{6}{{0}}%
\regfield{I2C\_TRANS\_START}{1}{5}{{0}}%
\regfield{I2C\_MS\_MODE}{1}{4}{{0}}%
\regfield{I2C\_RX\_FULL\_ACK\_LEVEL}{1}{3}{{1}}%
\regfield{I2C\_SAMPLE\_SCL\_LEVEL}{1}{2}{{0}}%
\regfield{I2C\_SCL\_FORCE\_OUT}{1}{1}{{1}}%
\regfield{I2C\_SDA\_FORCE\_OUT}{1}{0}{{1}}%
\reglabel{Reset}\regnewline%
\begin{regdesc}\begin{reglist}
\label{fielddesc:I2CSDAFORCEOUT}\item [I2C\_SDA\_FORCE\_OUT] 配置 SDA 输出模式。\\
0: 开漏输出 \\
1: 直接输出 \\ (读/写)
\label{fielddesc:I2CSCLFORCEOUT}\item [I2C\_SCL\_FORCE\_OUT] 配置 SCL 输出模式。\\
0: 开漏输出 \\
1: 直接输出 \\ (读/写)
\label{fielddesc:I2CSAMPLESCLLEVEL}\item [I2C\_SAMPLE\_SCL\_LEVEL] 用于选择采样模式。
1：SCL 为低电平时采样 SDA 数据。
0：SCL 为高电平时采样 SDA 数据。 (读/写)
\label{fielddesc:I2CRXFULLACKLEVEL}\item [I2C\_RX\_FULL\_ACK\_LEVEL] 用于配置主机在 rx\_fifo\_cnt 达到阈值时需发送的 ACK 电平值。 (读/写)
\label{fielddesc:I2CMSMODE}\item [I2C\_MS\_MODE] 置位此位，将模块配置为 I2C 主机。 清零此位，将模块配置为 I2C 从机。
 (读/写)
\label{fielddesc:I2CTRANSSTART}\item [I2C\_TRANS\_START] 置位此位，开始发送 TX FIFO 中的数据。 (读/写)
\label{fielddesc:I2CTXLSBFIRST}\item [I2C\_TX\_LSB\_FIRST] 用于控制待发送数据的发送模式。
1：从最低有效位开始发送数据；
0：从最高有效位开始发送数据。 (读/写)
\label{fielddesc:I2CRXLSBFIRST}\item [I2C\_RX\_LSB\_FIRST] 用于控制接收数据的存储模式。
1：从最低有效位开始接收数据；
0：从最高有效位开始接收数据。 (读/写)
\label{fielddesc:I2CCLKEN}\item [I2C\_CLK\_EN] 保留 (读/写)
\label{fielddesc:I2CARBITRATIONEN}\item [I2C\_ARBITRATION\_EN] I2C 总线仲裁的使能位。 (读/写)
\label{fielddesc:I2CFSMRST}\item [I2C\_FSM\_RST] 用于复位 SCL\_FSM。 (读/写)
\label{fielddesc:I2CREFALWAYSON}\item [I2C\_REF\_ALWAYS\_ON] 用于控制 REF\_TICK。 (读/写)
\end{reglist}\end{regdesc}
\end{register}


\begin{register}{H}{I2C\_TO\_REG}{\hexprefix{}000C}\label{regdesc:I2CTOREG}
\regfield{(reserved)}{7}{25}{0000000}%
\regfield{I2C\_TIME\_OUT\_EN}{1}{24}{{0}}%
\regfield{I2C\_TIME\_OUT\_VALUE}{24}{0}{{\hexprefix{}0000}}%
\reglabel{Reset}\regnewline%
\begin{regdesc}\begin{reglist}
\label{fielddesc:I2CTIMEOUTVALUE}\item [I2C\_TIME\_OUT\_VALUE] 用于配置接收一位数据的超时时间，以 APB
时钟周期为单位。 (读/写)
\label{fielddesc:I2CTIMEOUTEN}\item [I2C\_TIME\_OUT\_EN] 超时控制使能位。 (读/写)
\end{reglist}\end{regdesc}
\end{register}


\begin{register}{H}{I2C\_SLAVE\_ADDR\_REG}{\hexprefix{}0010}\label{regdesc:I2CSLAVEADDRREG}
\regfield{I2C\_ADDR\_10BIT\_EN}{1}{31}{{0}}%
\regfield{(reserved)}{16}{15}{0000000000000000}%
\regfield{I2C\_SLAVE\_ADDR}{15}{0}{{\hexprefix{}00}}%
\reglabel{Reset}\regnewline%
\begin{regdesc}\begin{reglist}
\label{fielddesc:I2CSLAVEADDR}\item [I2C\_SLAVE\_ADDR] 配置为 Slave 时，该字段用于配置从机地址。 (读/写)
\label{fielddesc:I2CADDR10BITEN}\item [I2C\_ADDR\_10BIT\_EN] 用于在主机模式下使能从机的 10 位寻址模式。 (读/写)
\end{reglist}\end{regdesc}
\end{register}


\begin{register}{H}{I2C\_FIFO\_CONF\_REG}{\hexprefix{}0018}\label{regdesc:I2CFIFOCONFREG}
\regfield{(reserved)}{5}{27}{00000}%
\regfield{I2C\_FIFO\_PRT\_EN}{1}{26}{{1}}%
\regfield{I2C\_NONFIFO\_TX\_THRES}{6}{20}{{\hexprefix{}15}}%
\regfield{I2C\_NONFIFO\_RX\_THRES}{6}{14}{{\hexprefix{}15}}%
\regfield{I2C\_TX\_FIFO\_RST}{1}{13}{{0}}%
\regfield{I2C\_RX\_FIFO\_RST}{1}{12}{{0}}%
\regfield{I2C\_FIFO\_ADDR\_CFG\_EN}{1}{11}{{0}}%
\regfield{I2C\_NONFIFO\_EN}{1}{10}{{0}}%
\regfield{I2C\_TXFIFO\_WM\_THRHD}{5}{5}{{\hexprefix{}4}}%
\regfield{I2C\_RXFIFO\_WM\_THRHD}{5}{0}{{\hexprefix{}b}}%
\reglabel{Reset}\regnewline%
\begin{regdesc}\begin{reglist}
\label{fielddesc:I2CRXFIFOWMTHRHD}\item [I2C\_RXFIFO\_WM\_THRHD] non-FIFO 访问模式下，RX FIFO 的水标阈值。I2C\_FIFO\_PRT\_EN 为 1 且 RX FIFO 计数值大于 I2C\_TXFIFO\_WM\_THRHD[4:0] 时，I2C\_TXFIFO\_WM\_INT\_RAW 位有效。 (读/写)
\label{fielddesc:I2CTXFIFOWMTHRHD}\item [I2C\_TXFIFO\_WM\_THRHD] non-FIFO 访问模式下，TX FIFO 的水标阈值。I2C\_FIFO\_PRT\_EN 为 1 且 TX FIFO 计数值小于 I2C\_TXFIFO\_WM\_THRHD[4:0] 时，I2C\_TXFIFO\_WM\_INT\_RAW 位有效。 (读/写)
\label{fielddesc:I2CNONFIFOEN}\item [I2C\_NONFIFO\_EN] 置位此位，使能 APB non-FIFO 访问。 (读/写)
\label{fielddesc:I2CFIFOADDRCFGEN}\item [I2C\_FIFO\_ADDR\_CFG\_EN] 此位置 1 时，从机接收地址字节的后一个字节为从机 RAM 中的偏移地址。 (读/写)
\label{fielddesc:I2CRXFIFORST}\item [I2C\_RX\_FIFO\_RST] 置位此位，复位 RX FIFO 。 (读/写)
\label{fielddesc:I2CTXFIFORST}\item [I2C\_TX\_FIFO\_RST] 置位此位，复位 TX FIFO 。 (读/写)
\label{fielddesc:I2CNONFIFORXTHRES}\item [I2C\_NONFIFO\_RX\_THRES] I2C 接收的数据超过 I2C\_NONFIFO\_TX\_THRES 字节时，生成 I2C\_RXFIFO\_UDF\_INT 中断，更新接收数据的当前偏移地址。 (读/写)
\label{fielddesc:I2CNONFIFOTXTHRES}\item [I2C\_NONFIFO\_TX\_THRES] I2C 发送的数据超过 I2C\_NONFIFO\_TX\_THRES 个字节时，
生成 I2C\_TXFIFO\_OVF\_INT 中断，更新发送数据的当前偏移地址。 (读/写)
\label{fielddesc:I2CFIFOPRTEN}\item [I2C\_FIFO\_PRT\_EN] non-FIFO 访问模式下 FIFO 指针的控制使能位。 该位控制 TX FIFO 和 RX FIFO 溢出、下溢、为满、为空时的有效位和中断。 (读/写)
\end{reglist}\end{regdesc}
\end{register}


\begin{register}{H}{I2C\_SCL\_SP\_CONF\_REG}{\hexprefix{}00A0}\label{regdesc:I2CSCLSPCONFREG}
\regfield{(reserved)}{24}{8}{000000000000000000000000}%
\regfield{I2C\_SDA\_PD\_EN}{1}{7}{{0}}%
\regfield{I2C\_SCL\_PD\_EN}{1}{6}{{0}}%
\regfield{I2C\_SCL\_RST\_SLV\_NUM}{5}{1}{{\hexprefix{}0}}%
\regfield{I2C\_SCL\_RST\_SLV\_EN}{1}{0}{{0}}%
\reglabel{Reset}\regnewline%
\begin{regdesc}\begin{reglist}
\label{fielddesc:I2CSCLRSTSLVEN}\item [I2C\_SCL\_RST\_SLV\_EN] I2C 主机处于空闲状态时，置位此位发送 SCL 脉冲。 脉冲数量为 I2C\_SCL\_RST\_SLV\_NUM[4:0]。 (读/写)
\label{fielddesc:I2CSCLRSTSLVNUM}\item [I2C\_SCL\_RST\_SLV\_NUM] 配置主机模式下生成的 SCL 脉冲。I2C\_SCL\_RST\_SLV\_EN 为 1 时有效。 (读/写)
\label{fielddesc:I2CSCLPDEN}\item [I2C\_SCL\_PD\_EN] 降低 I2C SCL 输出功耗的使能位。 1：不工作，降低功耗。 0：正常工作。 将 I2C\_SCL\_FORCE\_OUT 和 I2C\_SCL\_PD\_EN 置 1 延长 SCL 的低电平时间。 (读/写)
\label{fielddesc:I2CSDAPDEN}\item [I2C\_SDA\_PD\_EN] 降低 I2C SDA 输出功耗的使能位。 1：不工作，降低功耗。 0：正常工作。 将 I2C\_SDA\_FORCE\_OUT 和 I2C\_SDA\_PD\_EN 置 1 延长 SDA 的低电平时间。 (读/写)
\end{reglist}\end{regdesc}
\end{register}


\begin{register}{H}{I2C\_SCL\_STRETCH\_CONF\_REG}{\hexprefix{}00A4}\label{regdesc:I2CSCLSTRETCHCONFREG}
\regfield{(reserved)}{20}{12}{00000000000000000000}%
\regfield{I2C\_SLAVE\_SCL\_STRETCH\_CLR}{1}{11}{{0}}%
\regfield{I2C\_SLAVE\_SCL\_STRETCH\_EN}{1}{10}{{0}}%
\regfield{I2C\_STRETCH\_PROTECT\_NUM}{10}{0}{{\hexprefix{}0}}%
\reglabel{Reset}\regnewline%
\begin{regdesc}\begin{reglist}
\label{fielddesc:I2CSTRETCHPROTECTNUM}\item [I2C\_STRETCH\_PROTECT\_NUM] 配置 I2C 从机延长 SCL 低电平时间，以时钟周期为单位。 (读/写)
\label{fielddesc:I2CSLAVESCLSTRETCHEN}\item [I2C\_SLAVE\_SCL\_STRETCH\_EN] 从机 SCL 延展传输功能的使能位。 1：使能。 0：关闭。I2C\_SLAVE\_SCL\_STRETCH\_EN 为 1 时延展时钟，延长 SCL 输出线的低电平时间。 延展传输的原因可见 I2C\_STRETCH\_CAUSE。 (读/写)
\label{fielddesc:I2CSLAVESCLSTRETCHCLR}\item [I2C\_SLAVE\_SCL\_STRETCH\_CLR] 置位此位，清除 I2C 从机的 SCL 延展传输功能。 (只写)
\end{reglist}\end{regdesc}
\end{register}


\begin{register}{H}{I2C\_SR\_REG}{\hexprefix{}0008}\label{regdesc:I2CSRREG}
\regfield{(reserved)}{1}{31}{0}%
\regfield{I2C\_SCL\_STATE\_LAST}{3}{28}{{\hexprefix{}0}}%
\regfield{(reserved)}{1}{27}{0}%
\regfield{I2C\_SCL\_MAIN\_STATE\_LAST}{3}{24}{{\hexprefix{}0}}%
\regfield{I2C\_TXFIFO\_CNT}{6}{18}{{\hexprefix{}0}}%
\regfield{(reserved)}{2}{16}{00}%
\regfield{I2C\_STRETCH\_CAUSE}{2}{14}{{\hexprefix{}0}}%
\regfield{I2C\_RXFIFO\_CNT}{6}{8}{{\hexprefix{}0}}%
\regfield{(reserved)}{1}{7}{0}%
\regfield{I2C\_BYTE\_TRANS}{1}{6}{{0}}%
\regfield{I2C\_SLAVE\_ADDRESSED}{1}{5}{{0}}%
\regfield{I2C\_BUS\_BUSY}{1}{4}{{0}}%
\regfield{I2C\_ARB\_LOST}{1}{3}{{0}}%
\regfield{I2C\_TIME\_OUT}{1}{2}{{0}}%
\regfield{I2C\_SLAVE\_RW}{1}{1}{{0}}%
\regfield{I2C\_RESP\_REC}{1}{0}{{0}}%
\reglabel{Reset}\regnewline%
\begin{regdesc}\begin{reglist}
\label{fielddesc:I2CRESPREC}\item [I2C\_RESP\_REC] 主机模式或从机模式下接收的 ACK 电平值。 0：ACK，1：NACK。 (只读)
\label{fielddesc:I2CSLAVERW}\item [I2C\_SLAVE\_RW] 从机模式下，1：主机读取从机数据；0：主机向从机写入数据。 (只读)
\label{fielddesc:I2CTIMEOUT}\item [I2C\_TIME\_OUT] I2C 控制器接收一位数据的时间超过 I2C\_TIME\_OUT 周期时，
该字段变为 1。 (只读)
\label{fielddesc:I2CARBLOST}\item [I2C\_ARB\_LOST] I2C 控制器不控制 SCL 线时，该寄存器变为 1 。 (只读)
\label{fielddesc:I2CBUSBUSY}\item [I2C\_BUS\_BUSY] 1：I2C 总线正在传输数据；0：I2C 总线处于空闲状态。 (只读)
\label{fielddesc:I2CSLAVEADDRESSED}\item [I2C\_SLAVE\_ADDRESSED] 配置成 I2C 从机、且主机发送地址
与从机地址匹配时，该位翻转为高电平。 (只读)
\label{fielddesc:I2CBYTETRANS}\item [I2C\_BYTE\_TRANS] 传输一个字节后，该字段变为 1。 (只读)
\label{fielddesc:I2CRXFIFOCNT}\item [I2C\_RXFIFO\_CNT] 该字段为需发送数据的字节数。 (只读)
\label{fielddesc:I2CSTRETCHCAUSE}\item [I2C\_STRETCH\_CAUSE] 从机模式下延长 SCL 低电平时间的原因。 0： I2C 开始读取数据时延长 SCL 的低电平时间。 1：从机模式下 TX FIFO 为空时延长 SCL 的低电平时间。 2：从机模式下 RX FIFO 为满时延长 SCL 的低电平时间。 (只读)
\label{fielddesc:I2CTXFIFOCNT}\item [I2C\_TXFIFO\_CNT] 该字段存储 RAM 接收数据的字节数。 (只读)
\label{fielddesc:I2CSCLMAINSTATELAST}\item [I2C\_SCL\_MAIN\_STATE\_LAST] 该字段为 I2C 模块状态机的状态。
0：空闲；1：地址偏移；2：ACK 地址；3：接收数据；4：发送数据；5：发送 ACK；6：等待 ACK (只读)
\label{fielddesc:I2CSCLSTATELAST}\item [I2C\_SCL\_STATE\_LAST] 该字段为生成 SCL 的状态机状态。
0：空闲状态；1：开始；2：下降沿；3：低电平；4：上升沿；5：高电平；6：停止 (只读)
\end{reglist}\end{regdesc}
\end{register}


\begin{register}{H}{I2C\_FIFO\_ST\_REG}{\hexprefix{}0014}\label{regdesc:I2CFIFOSTREG}
\regfield{(reserved)}{2}{30}{00}%
\regfield{I2C\_SLAVE\_RW\_POINT}{8}{22}{{\hexprefix{}0}}%
\regfield{I2C\_TX\_UPDATE}{1}{21}{{0}}%
\regfield{I2C\_RX\_UPDATE}{1}{20}{{0}}%
\regfield{I2C\_TXFIFO\_END\_ADDR}{5}{15}{{\hexprefix{}0}}%
\regfield{I2C\_TXFIFO\_START\_ADDR}{5}{10}{{\hexprefix{}0}}%
\regfield{I2C\_RXFIFO\_END\_ADDR}{5}{5}{{\hexprefix{}0}}%
\regfield{I2C\_RXFIFO\_START\_ADDR}{5}{0}{{\hexprefix{}0}}%
\reglabel{Reset}\regnewline%
\begin{regdesc}\begin{reglist}
\label{fielddesc:I2CRXFIFOSTARTADDR}\item [I2C\_RXFIFO\_START\_ADDR] 最后接收数据的偏移地址，如寄存器 I2C\_NONFIFO\_RX\_THRES 所述。 (只读)
\label{fielddesc:I2CRXFIFOENDADDR}\item [I2C\_RXFIFO\_END\_ADDR] 最后接收数据的偏移地址，如寄存器 I2C\_NONFIFO\_RX\_THRES 所述。 该值在 I2C\_RX\_REC\_FULL\_INT 中断或 I2C\_TRANS\_COMPLETE\_INT 中断产生时更新。 (只读)
\label{fielddesc:I2CTXFIFOSTARTADDR}\item [I2C\_TXFIFO\_START\_ADDR] 最先发送数据的偏移地址，如寄存器 I2C\_NONFIFO\_TX\_THRES 所述。 (只读)
\label{fielddesc:I2CTXFIFOENDADDR}\item [I2C\_TXFIFO\_END\_ADDR] 最后发送数据的偏移地址，如寄存器 I2C\_NONFIFO\_TX\_THRES 所述。
该值在 I2C\_TX\_SEND\_EMPTY\_INT 中断或 I2C\_TRANS\_COMPLETE\_INT 中断产生时更新。 (只读)
\label{fielddesc:I2CRXUPDATE}\item [I2C\_RX\_UPDATE] 在 I2C\_RX\_UPDATE 上写 0 或 1，更新 I2C\_RXFIFO\_END\_ADDR 和 I2C\_RXFIFO\_START\_ADDR 的值。 (只写)
\label{fielddesc:I2CTXUPDATE}\item [I2C\_TX\_UPDATE] 在 I2C\_TX\_UPDATE 上写 0 或 1，更新 I2C\_TXFIFO\_END\_ADDR 和 I2C\_TXFIFO\_START\_ADDR 的值。 (只写)
\label{fielddesc:I2CSLAVERWPOINT}\item [I2C\_SLAVE\_RW\_POINT] 从机模式下接收的数据。 (只读)
\end{reglist}\end{regdesc}
\end{register}


\begin{register}{H}{I2C\_DATA\_REG}{\hexprefix{}001C}\label{regdesc:I2CDATAREG}
\regfield{(reserved)}{24}{8}{000000000000000000000000}%
\regfield{I2C\_FIFO\_RDATA}{8}{0}{{\hexprefix{}0}}%
\reglabel{Reset}\regnewline%
\begin{regdesc}\begin{reglist}
\label{fielddesc:I2CFIFORDATA}\item [I2C\_FIFO\_RDATA] 用于读取 RX FIFO 的数据，或向 TX FIFO 写数据。 (读/写)
\end{reglist}\end{regdesc}
\end{register}


\begin{register}{H}{I2C\_INT\_RAW\_REG}{\hexprefix{}0020}\label{regdesc:I2CINTRAWREG}
\regfield{(reserved)}{15}{17}{000000000000000}%
\regfield{I2C\_SLAVE\_STRETCH\_INT\_RAW}{1}{16}{{0}}%
\regfield{I2C\_DET\_START\_INT\_RAW}{1}{15}{{0}}%
\regfield{I2C\_SCL\_MAIN\_ST\_TO\_INT\_RAW}{1}{14}{{0}}%
\regfield{I2C\_SCL\_ST\_TO\_INT\_RAW}{1}{13}{{0}}%
\regfield{I2C\_RXFIFO\_UDF\_INT\_RAW}{1}{12}{{0}}%
\regfield{I2C\_TXFIFO\_OVF\_INT\_RAW}{1}{11}{{0}}%
\regfield{I2C\_NACK\_INT\_RAW}{1}{10}{{0}}%
\regfield{I2C\_TRANS\_START\_INT\_RAW}{1}{9}{{0}}%
\regfield{I2C\_TIME\_OUT\_INT\_RAW}{1}{8}{{0}}%
\regfield{I2C\_TRANS\_COMPLETE\_INT\_RAW}{1}{7}{{0}}%
\regfield{I2C\_MST\_TXFIFO\_UDF\_INT\_RAW}{1}{6}{{0}}%
\regfield{I2C\_ARBITRATION\_LOST\_INT\_RAW}{1}{5}{{0}}%
\regfield{I2C\_BYTE\_TRANS\_DONE\_INT\_RAW}{1}{4}{{0}}%
\regfield{I2C\_END\_DETECT\_INT\_RAW}{1}{3}{{0}}%
\regfield{I2C\_RXFIFO\_OVF\_INT\_RAW}{1}{2}{{0}}%
\regfield{I2C\_TXFIFO\_WM\_INT\_RAW}{1}{1}{{0}}%
\regfield{I2C\_RXFIFO\_WM\_INT\_RAW}{1}{0}{{0}}%
\reglabel{Reset}\regnewline%
\begin{regdesc}\begin{reglist}
\label{fielddesc:I2CRXFIFOWMINTRAW}\item [I2C\_RXFIFO\_WM\_INT\_RAW] I2C\_RXFIFO\_WM\_INT 的原始中断位。 (只读)
\label{fielddesc:I2CTXFIFOWMINTRAW}\item [I2C\_TXFIFO\_WM\_INT\_RAW] I2C\_TXFIFO\_WM\_INT 的原始中断位。 (只读)
\label{fielddesc:I2CRXFIFOOVFINTRAW}\item [I2C\_RXFIFO\_OVF\_INT\_RAW] I2C\_RXFIFO\_OVF\_INT 的原始中断位。 (只读)
\label{fielddesc:I2CENDDETECTINTRAW}\item [I2C\_END\_DETECT\_INT\_RAW] I2C\_END\_DETECT\_INT 的原始中断位。 (只读)
\label{fielddesc:I2CBYTETRANSDONEINTRAW}\item [I2C\_BYTE\_TRANS\_DONE\_INT\_RAW] I2C\_BYTE\_TRANS\_DONE\_INT 的原始中断位。 (只读)
\label{fielddesc:I2CARBITRATIONLOSTINTRAW}\item [I2C\_ARBITRATION\_LOST\_INT\_RAW] I2C\_ARBITRATION\_LOST\_INT 的原始中断位。 (只读)
\label{fielddesc:I2CMSTTXFIFOUDFINTRAW}\item [I2C\_MST\_TXFIFO\_UDF\_INT\_RAW] I2C\_MST\_TXFIFO\_UDF\_INT 的原始中断位。 (只读)
\label{fielddesc:I2CTRANSCOMPLETEINTRAW}\item [I2C\_TRANS\_COMPLETE\_INT\_RAW] I2C\_TRANS\_COMPLETE\_INT 的原始中断位。 (只读)
\label{fielddesc:I2CTIMEOUTINTRAW}\item [I2C\_TIME\_OUT\_INT\_RAW] I2C\_TIME\_OUT\_INT 的原始中断位。 (只读)
\label{fielddesc:I2CTRANSSTARTINTRAW}\item [I2C\_TRANS\_START\_INT\_RAW] I2C\_TRANS\_START\_INT 的原始中断位。 (只读)
\label{fielddesc:I2CNACKINTRAW}\item [I2C\_NACK\_INT\_RAW] I2C\_NACK\_INT 的原始中断位。 (只读)
\label{fielddesc:I2CTXFIFOOVFINTRAW}\item [I2C\_TXFIFO\_OVF\_INT\_RAW] I2C\_TXFIFO\_OVF\_INT 的原始中断位。 (只读)
\label{fielddesc:I2CRXFIFOUDFINTRAW}\item [I2C\_RXFIFO\_UDF\_INT\_RAW] I2C\_RXFIFO\_UDF\_INT 的原始中断位。 (只读)
\label{fielddesc:I2CSCLSTTOINTRAW}\item [I2C\_SCL\_ST\_TO\_INT\_RAW] I2C\_SCL\_ST\_TO\_INT 的原始中断位。 (只读)
\label{fielddesc:I2CSCLMAINSTTOINTRAW}\item [I2C\_SCL\_MAIN\_ST\_TO\_INT\_RAW] I2C\_SCL\_MAIN\_ST\_TO\_INT 的原始中断位。 (只读)
\label{fielddesc:I2CDETSTARTINTRAW}\item [I2C\_DET\_START\_INT\_RAW] I2C\_DET\_START\_INT 的原始中断位。 (只读)
\label{fielddesc:I2CSLAVESTRETCHINTRAW}\item [I2C\_SLAVE\_STRETCH\_INT\_RAW] I2C\_SLAVE\_STRETCH\_INT 的原始中断位。 (只读)
\end{reglist}\end{regdesc}
\end{register}


\begin{register}{H}{I2C\_INT\_CLR\_REG}{\hexprefix{}0024}\label{regdesc:I2CINTCLRREG}
\regfield{(reserved)}{15}{17}{000000000000000}%
\regfield{I2C\_SLAVE\_STRETCH\_INT\_CLR}{1}{16}{{0}}%
\regfield{I2C\_DET\_START\_INT\_CLR}{1}{15}{{0}}%
\regfield{I2C\_SCL\_MAIN\_ST\_TO\_INT\_CLR}{1}{14}{{0}}%
\regfield{I2C\_SCL\_ST\_TO\_INT\_CLR}{1}{13}{{0}}%
\regfield{I2C\_RXFIFO\_UDF\_INT\_CLR}{1}{12}{{0}}%
\regfield{I2C\_TXFIFO\_OVF\_INT\_CLR}{1}{11}{{0}}%
\regfield{I2C\_NACK\_INT\_CLR}{1}{10}{{0}}%
\regfield{I2C\_TRANS\_START\_INT\_CLR}{1}{9}{{0}}%
\regfield{I2C\_TIME\_OUT\_INT\_CLR}{1}{8}{{0}}%
\regfield{I2C\_TRANS\_COMPLETE\_INT\_CLR}{1}{7}{{0}}%
\regfield{I2C\_MST\_TXFIFO\_UDF\_INT\_CLR}{1}{6}{{0}}%
\regfield{I2C\_ARBITRATION\_LOST\_INT\_CLR}{1}{5}{{0}}%
\regfield{I2C\_BYTE\_TRANS\_DONE\_INT\_CLR}{1}{4}{{0}}%
\regfield{I2C\_END\_DETECT\_INT\_CLR}{1}{3}{{0}}%
\regfield{I2C\_RXFIFO\_OVF\_INT\_CLR}{1}{2}{{0}}%
\regfield{I2C\_TXFIFO\_WM\_INT\_CLR}{1}{1}{{0}}%
\regfield{I2C\_RXFIFO\_WM\_INT\_CLR}{1}{0}{{0}}%
\reglabel{Reset}\regnewline%
\begin{regdesc}\begin{reglist}
\label{fielddesc:I2CRXFIFOWMINTCLR}\item [I2C\_RXFIFO\_WM\_INT\_CLR] 置位此位，清除 I2C\_RXFIFO\_WM\_INT 中断。 (只写)
\label{fielddesc:I2CTXFIFOWMINTCLR}\item [I2C\_TXFIFO\_WM\_INT\_CLR] 置位此位，清除 I2C\_TXFIFO\_WM\_INT 中断。 (只写)
\label{fielddesc:I2CRXFIFOOVFINTCLR}\item [I2C\_RXFIFO\_OVF\_INT\_CLR] 置位此位，清除 I2C\_RXFIFO\_OVF\_INT 中断。 (只写)
\label{fielddesc:I2CENDDETECTINTCLR}\item [I2C\_END\_DETECT\_INT\_CLR] 置位此位，清除 I2C\_END\_DETECT\_INT 中断。 (只写)
\label{fielddesc:I2CBYTETRANSDONEINTCLR}\item [I2C\_BYTE\_TRANS\_DONE\_INT\_CLR] 置位此位，清除 I2C\_BYTE\_TRANS\_DONE\_INT 中断。 (只写)
\label{fielddesc:I2CARBITRATIONLOSTINTCLR}\item [I2C\_ARBITRATION\_LOST\_INT\_CLR] 置位此位，清除 I2C\_ARBITRATION\_LOST\_INT 中断。 (只写)
\label{fielddesc:I2CMSTTXFIFOUDFINTCLR}\item [I2C\_MST\_TXFIFO\_UDF\_INT\_CLR] 置位此位，清除 I2C\_MST\_TXFIFO\_UDF\_INT 中断。 (只写)
\label{fielddesc:I2CTRANSCOMPLETEINTCLR}\item [I2C\_TRANS\_COMPLETE\_INT\_CLR] 置位此位，清除 I2C\_TRANS\_COMPLETE\_INT 中断。 (只写)
\label{fielddesc:I2CTIMEOUTINTCLR}\item [I2C\_TIME\_OUT\_INT\_CLR] 置位此位，清除 I2C\_TIME\_OUT\_INT 中断。 (只写)
\label{fielddesc:I2CTRANSSTARTINTCLR}\item [I2C\_TRANS\_START\_INT\_CLR] 置位此位，清除 I2C\_TRANS\_START\_INT 中断。 (只写)
\label{fielddesc:I2CNACKINTCLR}\item [I2C\_NACK\_INT\_CLR] 置位此位，清除 I2C\_NACK\_INT 中断。 (只写)
\label{fielddesc:I2CTXFIFOOVFINTCLR}\item [I2C\_TXFIFO\_OVF\_INT\_CLR] 置位此位，清除 I2C\_TXFIFO\_OVF\_INT 中断。 (只写)
\label{fielddesc:I2CRXFIFOUDFINTCLR}\item [I2C\_RXFIFO\_UDF\_INT\_CLR] 置位此位，清除 I2C\_RXFIFO\_UDF\_INT 中断。 (只写)
\label{fielddesc:I2CSCLSTTOINTCLR}\item [I2C\_SCL\_ST\_TO\_INT\_CLR] 置位此位，清除 I2C\_SCL\_ST\_TO\_INT 中断。 (只写)
\label{fielddesc:I2CSCLMAINSTTOINTCLR}\item [I2C\_SCL\_MAIN\_ST\_TO\_INT\_CLR] 置位此位，清除 I2C\_SCL\_MAIN\_ST\_TO\_INT 中断。 (只写)
\label{fielddesc:I2CDETSTARTINTCLR}\item [I2C\_DET\_START\_INT\_CLR] 置位此位，清除 I2C\_DET\_START\_INT 中断。 (只写)
\label{fielddesc:I2CSLAVESTRETCHINTCLR}\item [I2C\_SLAVE\_STRETCH\_INT\_CLR] 置位此位，清除 I2C\_SLAVE\_STRETCH\_INT 中断。 (只写)
\end{reglist}\end{regdesc}
\end{register}


\begin{register}{H}{I2C\_INT\_ENA\_REG}{\hexprefix{}0028}\label{regdesc:I2CINTENAREG}
\regfield{(reserved)}{15}{17}{000000000000000}%
\regfield{I2C\_SLAVE\_STRETCH\_INT\_ENA}{1}{16}{{0}}%
\regfield{I2C\_DET\_START\_INT\_ENA}{1}{15}{{0}}%
\regfield{I2C\_SCL\_MAIN\_ST\_TO\_INT\_ENA}{1}{14}{{0}}%
\regfield{I2C\_SCL\_ST\_TO\_INT\_ENA}{1}{13}{{0}}%
\regfield{I2C\_RXFIFO\_UDF\_INT\_ENA}{1}{12}{{0}}%
\regfield{I2C\_TXFIFO\_OVF\_INT\_ENA}{1}{11}{{0}}%
\regfield{I2C\_NACK\_INT\_ENA}{1}{10}{{0}}%
\regfield{I2C\_TRANS\_START\_INT\_ENA}{1}{9}{{0}}%
\regfield{I2C\_TIME\_OUT\_INT\_ENA}{1}{8}{{0}}%
\regfield{I2C\_TRANS\_COMPLETE\_INT\_ENA}{1}{7}{{0}}%
\regfield{I2C\_MST\_TXFIFO\_UDF\_INT\_ENA}{1}{6}{{0}}%
\regfield{I2C\_ARBITRATION\_LOST\_INT\_ENA}{1}{5}{{0}}%
\regfield{I2C\_BYTE\_TRANS\_DONE\_INT\_ENA}{1}{4}{{0}}%
\regfield{I2C\_END\_DETECT\_INT\_ENA}{1}{3}{{0}}%
\regfield{I2C\_RXFIFO\_OVF\_INT\_ENA}{1}{2}{{0}}%
\regfield{I2C\_TXFIFO\_WM\_INT\_ENA}{1}{1}{{0}}%
\regfield{I2C\_RXFIFO\_WM\_INT\_ENA}{1}{0}{{0}}%
\reglabel{Reset}\regnewline%
\begin{regdesc}\begin{reglist}
\label{fielddesc:I2CRXFIFOWMINTENA}\item [I2C\_RXFIFO\_WM\_INT\_ENA] I2C\_RXFIFO\_WM\_INT 的原始中断位。 (读/写)
\label{fielddesc:I2CTXFIFOWMINTENA}\item [I2C\_TXFIFO\_WM\_INT\_ENA] I2C\_TXFIFO\_WM\_INT 的原始中断位。 (读/写)
\label{fielddesc:I2CRXFIFOOVFINTENA}\item [I2C\_RXFIFO\_OVF\_INT\_ENA] I2C\_RXFIFO\_OVF\_INT 的原始中断位。 (读/写)
\label{fielddesc:I2CENDDETECTINTENA}\item [I2C\_END\_DETECT\_INT\_ENA] I2C\_END\_DETECT\_INT 的原始中断位。 (读/写)
\label{fielddesc:I2CBYTETRANSDONEINTENA}\item [I2C\_BYTE\_TRANS\_DONE\_INT\_ENA] I2C\_BYTE\_TRANS\_DONE\_INT 的原始中断位。 (读/写)
\label{fielddesc:I2CARBITRATIONLOSTINTENA}\item [I2C\_ARBITRATION\_LOST\_INT\_ENA] I2C\_ARBITRATION\_LOST\_INT 的原始中断位。 (读/写)
\label{fielddesc:I2CMSTTXFIFOUDFINTENA}\item [I2C\_MST\_TXFIFO\_UDF\_INT\_ENA] I2C\_MST\_TXFIFO\_UDF\_INT 的原始中断位。 (读/写)
\label{fielddesc:I2CTRANSCOMPLETEINTENA}\item [I2C\_TRANS\_COMPLETE\_INT\_ENA] I2C\_TRANS\_COMPLETE\_INT 的原始中断位。 (读/写)
\label{fielddesc:I2CTIMEOUTINTENA}\item [I2C\_TIME\_OUT\_INT\_ENA] I2C\_TIME\_OUT\_INT 的原始中断位。 (读/写)
\label{fielddesc:I2CTRANSSTARTINTENA}\item [I2C\_TRANS\_START\_INT\_ENA] I2C\_TRANS\_START\_INT 的原始中断位。 (读/写)
\label{fielddesc:I2CNACKINTENA}\item [I2C\_NACK\_INT\_ENA] I2C\_NACK\_INT 的原始中断位。 (读/写)
\label{fielddesc:I2CTXFIFOOVFINTENA}\item [I2C\_TXFIFO\_OVF\_INT\_ENA] I2C\_TXFIFO\_OVF\_INT 的原始中断位。 (读/写)
\label{fielddesc:I2CRXFIFOUDFINTENA}\item [I2C\_RXFIFO\_UDF\_INT\_ENA] I2C\_RXFIFO\_UDF\_INT 的原始中断位。 (读/写)
\label{fielddesc:I2CSCLSTTOINTENA}\item [I2C\_SCL\_ST\_TO\_INT\_ENA] I2C\_SCL\_ST\_TO\_INT 的原始中断位。 (读/写)
\label{fielddesc:I2CSCLMAINSTTOINTENA}\item [I2C\_SCL\_MAIN\_ST\_TO\_INT\_ENA] I2C\_SCL\_MAIN\_ST\_TO\_INT 的原始中断位。 (读/写)
\label{fielddesc:I2CDETSTARTINTENA}\item [I2C\_DET\_START\_INT\_ENA] I2C\_DET\_START\_INT 的原始中断位。 (读/写)
\label{fielddesc:I2CSLAVESTRETCHINTENA}\item [I2C\_SLAVE\_STRETCH\_INT\_ENA] I2C\_SLAVE\_STRETCH\_INT 的原始中断位。 (读/写)
\end{reglist}\end{regdesc}
\end{register}


\begin{register}{H}{I2C\_INT\_STATUS\_REG}{\hexprefix{}002C}\label{regdesc:I2CINTSTATUSREG}
\regfield{(reserved)}{15}{17}{000000000000000}%
\regfield{I2C\_SLAVE\_STRETCH\_INT\_ST}{1}{16}{{0}}%
\regfield{I2C\_DET\_START\_INT\_ST}{1}{15}{{0}}%
\regfield{I2C\_SCL\_MAIN\_ST\_TO\_INT\_ST}{1}{14}{{0}}%
\regfield{I2C\_SCL\_ST\_TO\_INT\_ST}{1}{13}{{0}}%
\regfield{I2C\_RXFIFO\_UDF\_INT\_ST}{1}{12}{{0}}%
\regfield{I2C\_TXFIFO\_OVF\_INT\_ST}{1}{11}{{0}}%
\regfield{I2C\_NACK\_INT\_ST}{1}{10}{{0}}%
\regfield{I2C\_TRANS\_START\_INT\_ST}{1}{9}{{0}}%
\regfield{I2C\_TIME\_OUT\_INT\_ST}{1}{8}{{0}}%
\regfield{I2C\_TRANS\_COMPLETE\_INT\_ST}{1}{7}{{0}}%
\regfield{I2C\_MST\_TXFIFO\_UDF\_INT\_ST}{1}{6}{{0}}%
\regfield{I2C\_ARBITRATION\_LOST\_INT\_ST}{1}{5}{{0}}%
\regfield{I2C\_BYTE\_TRANS\_DONE\_INT\_ST}{1}{4}{{0}}%
\regfield{I2C\_END\_DETECT\_INT\_ST}{1}{3}{{0}}%
\regfield{I2C\_RXFIFO\_OVF\_INT\_ST}{1}{2}{{0}}%
\regfield{I2C\_TXFIFO\_WM\_INT\_ST}{1}{1}{{0}}%
\regfield{I2C\_RXFIFO\_WM\_INT\_ST}{1}{0}{{0}}%
\reglabel{Reset}\regnewline%
\begin{regdesc}\begin{reglist}
\label{fielddesc:I2CRXFIFOWMINTST}\item [I2C\_RXFIFO\_WM\_INT\_ST] I2C\_RXFIFO\_WM\_INT 的屏蔽中断状态位。 (只读)
\label{fielddesc:I2CTXFIFOWMINTST}\item [I2C\_TXFIFO\_WM\_INT\_ST] I2C\_TXFIFO\_WM\_INT 的屏蔽中断状态位。 (只读)
\label{fielddesc:I2CRXFIFOOVFINTST}\item [I2C\_RXFIFO\_OVF\_INT\_ST] I2C\_RXFIFO\_OVF\_INT 的屏蔽中断状态位。 (只读)
\label{fielddesc:I2CENDDETECTINTST}\item [I2C\_END\_DETECT\_INT\_ST] I2C\_END\_DETECT\_INT 的屏蔽中断状态位。 (只读)
\label{fielddesc:I2CBYTETRANSDONEINTST}\item [I2C\_BYTE\_TRANS\_DONE\_INT\_ST] I2C\_BYTE\_TRANS\_DONE\_INT 的屏蔽中断状态位。 (只读)
\label{fielddesc:I2CARBITRATIONLOSTINTST}\item [I2C\_ARBITRATION\_LOST\_INT\_ST] I2C\_ARBITRATION\_LOST\_INT 的屏蔽中断状态位。 (只读)
\label{fielddesc:I2CMSTTXFIFOUDFINTST}\item [I2C\_MST\_TXFIFO\_UDF\_INT\_ST] I2C\_MST\_TXFIFO\_UDF\_INT 的屏蔽中断状态位。 (只读)
\label{fielddesc:I2CTRANSCOMPLETEINTST}\item [I2C\_TRANS\_COMPLETE\_INT\_ST] I2C\_TRANS\_COMPLETE\_INT 的屏蔽中断状态位。 (只读)
\label{fielddesc:I2CTIMEOUTINTST}\item [I2C\_TIME\_OUT\_INT\_ST] I2C\_TIME\_OUT\_INT 的屏蔽中断状态位。 (只读)
\label{fielddesc:I2CTRANSSTARTINTST}\item [I2C\_TRANS\_START\_INT\_ST] I2C\_TRANS\_START\_INT 的屏蔽中断状态位。 (只读)
\label{fielddesc:I2CNACKINTST}\item [I2C\_NACK\_INT\_ST] I2C\_NACK\_INT 的屏蔽中断状态位。 (只读)
\label{fielddesc:I2CTXFIFOOVFINTST}\item [I2C\_TXFIFO\_OVF\_INT\_ST] I2C\_TXFIFO\_OVF\_INT 的屏蔽中断状态位。 (只读)
\label{fielddesc:I2CRXFIFOUDFINTST}\item [I2C\_RXFIFO\_UDF\_INT\_ST] I2C\_RXFIFO\_UDF\_INT 的屏蔽中断状态位。 (只读)
\label{fielddesc:I2CSCLSTTOINTST}\item [I2C\_SCL\_ST\_TO\_INT\_ST] I2C\_SCL\_ST\_TO\_INT 的屏蔽中断状态位。 (只读)
\label{fielddesc:I2CSCLMAINSTTOINTST}\item [I2C\_SCL\_MAIN\_ST\_TO\_INT\_ST] I2C\_SCL\_MAIN\_ST\_TO\_INT 的屏蔽中断状态位。 (只读)
\label{fielddesc:I2CDETSTARTINTST}\item [I2C\_DET\_START\_INT\_ST] I2C\_DET\_START\_INT 的屏蔽中断状态位。 (只读)
\label{fielddesc:I2CSLAVESTRETCHINTST}\item [I2C\_SLAVE\_STRETCH\_INT\_ST] I2C\_SLAVE\_STRETCH\_INT 的屏蔽中断状态位。 (只读)
\end{reglist}\end{regdesc}
\end{register}


\begin{register}{H}{I2C\_SCL\_FILTER\_CFG\_REG}{\hexprefix{}0050}\label{regdesc:I2CSCLFILTERCFGREG}
\regfield{(reserved)}{27}{5}{000000000000000000000000000}%
\regfield{I2C\_SCL\_FILTER\_EN}{1}{4}{{1}}%
\regfield{I2C\_SCL\_FILTER\_THRES}{4}{0}{{\hexprefix{}0}}%
\reglabel{Reset}\regnewline%
\begin{regdesc}\begin{reglist}
\label{fielddesc:I2CSCLFILTERTHRES}\item [I2C\_SCL\_FILTER\_THRES] SCL 输入信号的脉冲宽度小于该寄存器
的值时，I2C 控制器忽略此脉冲。该寄存器的值以 I2C 模块时钟周期数为单位。 (读/写)
\label{fielddesc:I2CSCLFILTEREN}\item [I2C\_SCL\_FILTER\_EN] SCL 的滤波使能位。 (读/写)
\end{reglist}\end{regdesc}
\end{register}


\begin{register}{H}{I2C\_SDA\_FILTER\_CFG\_REG}{\hexprefix{}0054}\label{regdesc:I2CSDAFILTERCFGREG}
\regfield{(reserved)}{27}{5}{000000000000000000000000000}%
\regfield{I2C\_SDA\_FILTER\_EN}{1}{4}{{1}}%
\regfield{I2C\_SDA\_FILTER\_THRES}{4}{0}{{\hexprefix{}0}}%
\reglabel{Reset}\regnewline%
\begin{regdesc}\begin{reglist}
\label{fielddesc:I2CSDAFILTERTHRES}\item [I2C\_SDA\_FILTER\_THRES] SDA 输入信号的脉冲宽度小于该寄存器
的值时，I2C 控制器忽略此脉冲。该寄存器的值以 I2C 模块时钟周期数为单位。 (读/写)
\label{fielddesc:I2CSDAFILTEREN}\item [I2C\_SDA\_FILTER\_EN] SDA 的滤波使能位。 (读/写)
\end{reglist}\end{regdesc}
\end{register}


\begin{register}{H}{I2C\_COMD0\_REG}{\hexprefix{}0058}\label{regdesc:I2CCOMD0REG}
\regfield{I2C\_COMMAND0\_DONE}{1}{31}{{0}}%
\regfield{(reserved)}{17}{14}{00000000000000000}%
\regfield{I2C\_COMMAND0}{14}{0}{{\hexprefix{}00}}%
\reglabel{Reset}\regnewline%
\begin{regdesc}\begin{reglist}
\label{fielddesc:I2CCOMMAND0}\item [I2C\_COMMAND0] 命令 0 的内容。 该命令包括三个部分：
op\_code 为命令，0：RSTART；1：WRITE；2：READ；3：STOP；4：END。
byte\_num 表示需发送或接收的字节数。
ack\_check\_en、ack\_exp 和 ack 用于控制 ACK 位。 参阅 I2C cmd 结构获取更多
信息。 (读/写)
\label{fielddesc:I2CCOMMAND0DONE}\item [I2C\_COMMAND0\_DONE] 在 I2C 主机模式下完成命令 0 时，该位翻转为高
电平。 (读/写)
\end{reglist}\end{regdesc}
\end{register}


\begin{register}{H}{I2C\_COMD1\_REG}{\hexprefix{}005C}\label{regdesc:I2CCOMD1REG}
\regfield{I2C\_COMMAND1\_DONE}{1}{31}{{0}}%
\regfield{(reserved)}{17}{14}{00000000000000000}%
\regfield{I2C\_COMMAND1}{14}{0}{{\hexprefix{}00}}%
\reglabel{Reset}\regnewline%
\begin{regdesc}\begin{reglist}
\label{fielddesc:I2CCOMMAND1}\item [I2C\_COMMAND1] 命令 1 的内容。 该命令包括三个部分：
op\_code 为命令，0：RSTART；1：WRITE；2：READ；3：STOP；4：END。
byte\_num 表示需发送或接收的字节数。
ack\_check\_en、ack\_exp 和 ack 用于控制 ACK 位。 参阅 I2C cmd 结构获取更多
信息。 (读/写)
\label{fielddesc:I2CCOMMAND1DONE}\item [I2C\_COMMAND1\_DONE] 在 I2C 主机模式下完成命令 1 时，该位翻转为高
电平。 (读/写)
\end{reglist}\end{regdesc}
\end{register}


\begin{register}{H}{I2C\_COMD2\_REG}{\hexprefix{}0060}\label{regdesc:I2CCOMD2REG}
\regfield{I2C\_COMMAND2\_DONE}{1}{31}{{0}}%
\regfield{(reserved)}{17}{14}{00000000000000000}%
\regfield{I2C\_COMMAND2}{14}{0}{{\hexprefix{}00}}%
\reglabel{Reset}\regnewline%
\begin{regdesc}\begin{reglist}
\label{fielddesc:I2CCOMMAND2}\item [I2C\_COMMAND2] 命令 2 的内容。 该命令包括三个部分：
op\_code 为命令，0：RSTART；1：WRITE；2：READ；3：STOP；4：END。
byte\_num 表示需发送或接收的字节数。
ack\_check\_en、ack\_exp 和 ack 用于控制 ACK 位。 参阅 I2C cmd 结构获取更多
信息。 (读/写)
\label{fielddesc:I2CCOMMAND2DONE}\item [I2C\_COMMAND2\_DONE] 在 I2C 主机模式下完成命令 2 时，该位翻转为高
电平。 (读/写)
\end{reglist}\end{regdesc}
\end{register}


\begin{register}{H}{I2C\_COMD3\_REG}{\hexprefix{}0064}\label{regdesc:I2CCOMD3REG}
\regfield{I2C\_COMMAND3\_DONE}{1}{31}{{0}}%
\regfield{(reserved)}{17}{14}{00000000000000000}%
\regfield{I2C\_COMMAND3}{14}{0}{{\hexprefix{}00}}%
\reglabel{Reset}\regnewline%
\begin{regdesc}\begin{reglist}
\label{fielddesc:I2CCOMMAND3}\item [I2C\_COMMAND3] 命令 3 的内容。 该命令包括三个部分：
op\_code 为命令，0：RSTART；1：WRITE；2：READ；3：STOP；4：END。
byte\_num 表示需发送或接收的字节数。
ack\_check\_en、ack\_exp 和 ack 用于控制 ACK 位。 参阅 I2C cmd 结构获取更多
信息。 (读/写)
\label{fielddesc:I2CCOMMAND3DONE}\item [I2C\_COMMAND3\_DONE] 在 I2C 主机模式下完成命令 3 时，该位翻转为高
电平。 (读/写)
\end{reglist}\end{regdesc}
\end{register}


\begin{register}{H}{I2C\_COMD4\_REG}{\hexprefix{}0068}\label{regdesc:I2CCOMD4REG}
\regfield{I2C\_COMMAND4\_DONE}{1}{31}{{0}}%
\regfield{(reserved)}{17}{14}{00000000000000000}%
\regfield{I2C\_COMMAND4}{14}{0}{{\hexprefix{}00}}%
\reglabel{Reset}\regnewline%
\begin{regdesc}\begin{reglist}
\label{fielddesc:I2CCOMMAND4}\item [I2C\_COMMAND4] 命令 4 的内容。 该命令包括三个部分：
op\_code 为命令，0：RSTART；1：WRITE；2：READ；3：STOP；4：END。
byte\_num 表示需发送或接收的字节数。
ack\_check\_en、ack\_exp 和 ack 用于控制 ACK 位。 参阅 I2C cmd 结构获取更多
信息。 (读/写)
\label{fielddesc:I2CCOMMAND4DONE}\item [I2C\_COMMAND4\_DONE] 在 I2C 主机模式下完成命令 4 时，该位翻转为高
电平。 (读/写)
\end{reglist}\end{regdesc}
\end{register}


\begin{register}{H}{I2C\_COMD5\_REG}{\hexprefix{}006C}\label{regdesc:I2CCOMD5REG}
\regfield{I2C\_COMMAND5\_DONE}{1}{31}{{0}}%
\regfield{(reserved)}{17}{14}{00000000000000000}%
\regfield{I2C\_COMMAND5}{14}{0}{{\hexprefix{}00}}%
\reglabel{Reset}\regnewline%
\begin{regdesc}\begin{reglist}
\label{fielddesc:I2CCOMMAND5}\item [I2C\_COMMAND5] 命令 5 的内容。 该命令包括三个部分：
op\_code 为命令，0：RSTART；1：WRITE；2：READ；3：STOP；4：END。
byte\_num 表示需发送或接收的字节数。
ack\_check\_en、ack\_exp 和 ack 用于控制 ACK 位。 参阅 I2C cmd 结构获取更多
信息。 (读/写)
\label{fielddesc:I2CCOMMAND5DONE}\item [I2C\_COMMAND5\_DONE] 在 I2C 主机模式下完成命令 5 时，该位翻转为高电平。 (读/写)
\end{reglist}\end{regdesc}
\end{register}


\begin{register}{H}{I2C\_COMD6\_REG}{\hexprefix{}0070}\label{regdesc:I2CCOMD6REG}
\regfield{I2C\_COMMAND6\_DONE}{1}{31}{{0}}%
\regfield{(reserved)}{17}{14}{00000000000000000}%
\regfield{I2C\_COMMAND6}{14}{0}{{\hexprefix{}00}}%
\reglabel{Reset}\regnewline%
\begin{regdesc}\begin{reglist}
\label{fielddesc:I2CCOMMAND6}\item [I2C\_COMMAND6] 命令 6 的内容。 该命令包括三个部分：
op\_code 为命令，0：RSTART；1：WRITE；2：READ；3：STOP；4：END。
byte\_num 表示需发送或接收的字节数。
ack\_check\_en、ack\_exp 和 ack 用于控制 ACK 位。 参阅 I2C cmd 结构获取更多
信息。 (读/写)
\label{fielddesc:I2CCOMMAND6DONE}\item [I2C\_COMMAND6\_DONE] 在 I2C 主机模式下完成命令 6 时，该位翻转为高电平。 (读/写)
\end{reglist}\end{regdesc}
\end{register}


\begin{register}{H}{I2C\_COMD7\_REG}{\hexprefix{}0074}\label{regdesc:I2CCOMD7REG}
\regfield{I2C\_COMMAND7\_DONE}{1}{31}{{0}}%
\regfield{(reserved)}{17}{14}{00000000000000000}%
\regfield{I2C\_COMMAND7}{14}{0}{{\hexprefix{}00}}%
\reglabel{Reset}\regnewline%
\begin{regdesc}\begin{reglist}
\label{fielddesc:I2CCOMMAND7}\item [I2C\_COMMAND7] 命令 7 的内容。 该命令包括三个部分：
op\_code 为命令，0：RSTART；1：WRITE；2：READ；3：STOP；4：END。
byte\_num 表示需发送或接收的字节数。
ack\_check\_en、ack\_exp 和 ack 用于控制 ACK 位。 参阅 I2C cmd 结构获取更多
信息。 (读/写)
\label{fielddesc:I2CCOMMAND7DONE}\item [I2C\_COMMAND7\_DONE] 在 I2C 主机模式下完成命令 7 时，该位翻转为高电平。 (读/写)
\end{reglist}\end{regdesc}
\end{register}


\begin{register}{H}{I2C\_COMD8\_REG}{\hexprefix{}0078}\label{regdesc:I2CCOMD8REG}
\regfield{I2C\_COMMAND8\_DONE}{1}{31}{{0}}%
\regfield{(reserved)}{17}{14}{00000000000000000}%
\regfield{I2C\_COMMAND8}{14}{0}{{\hexprefix{}00}}%
\reglabel{Reset}\regnewline%
\begin{regdesc}\begin{reglist}
\label{fielddesc:I2CCOMMAND8}\item [I2C\_COMMAND8] 命令 8 的内容。 该命令包括三个部分：
op\_code 为命令，0：RSTART；1：WRITE；2：READ；3：STOP；4：END。
byte\_num 表示需发送或接收的字节数。
ack\_check\_en、ack\_exp 和 ack 用于控制 ACK 位。 参阅 I2C cmd 结构获取更多
信息。 (读/写)
\label{fielddesc:I2CCOMMAND8DONE}\item [I2C\_COMMAND8\_DONE] 在 I2C 主机模式下完成命令 8 时，该位翻转为高电平。 (读/写)
\end{reglist}\end{regdesc}
\end{register}


\begin{register}{H}{I2C\_COMD9\_REG}{\hexprefix{}007C}\label{regdesc:I2CCOMD9REG}
\regfield{I2C\_COMMAND9\_DONE}{1}{31}{{0}}%
\regfield{(reserved)}{17}{14}{00000000000000000}%
\regfield{I2C\_COMMAND9}{14}{0}{{\hexprefix{}00}}%
\reglabel{Reset}\regnewline%
\begin{regdesc}\begin{reglist}
\label{fielddesc:I2CCOMMAND9}\item [I2C\_COMMAND9] 命令 9 的内容。 该命令包括三个部分：
op\_code 为命令，0：RSTART；1：WRITE；2：READ；3：STOP；4：END。
byte\_num 表示需发送或接收的字节数。
ack\_check\_en、ack\_exp 和 ack 用于控制 ACK 位。 参阅 I2C cmd 结构获取更多
信息。 (读/写)
\label{fielddesc:I2CCOMMAND9DONE}\item [I2C\_COMMAND9\_DONE] 在 I2C 主机模式下完成命令 9 时，该位翻转为高电平。 (读/写)
\end{reglist}\end{regdesc}
\end{register}


\begin{register}{H}{I2C\_COMD10\_REG}{\hexprefix{}0080}\label{regdesc:I2CCOMD10REG}
\regfield{I2C\_COMMAND10\_DONE}{1}{31}{{0}}%
\regfield{(reserved)}{17}{14}{00000000000000000}%
\regfield{I2C\_COMMAND10}{14}{0}{{\hexprefix{}00}}%
\reglabel{Reset}\regnewline%
\begin{regdesc}\begin{reglist}
\label{fielddesc:I2CCOMMAND10}\item [I2C\_COMMAND10] 命令 10 的内容。 该命令包括三个部分：
op\_code 为命令，0：RSTART；1：WRITE；2：READ；3：STOP；4：END。
byte\_num 表示需发送或接收的字节数。
ack\_check\_en、ack\_exp 和 ack 用于控制 ACK 位。 参阅 I2C cmd 结构获取更多
信息。 (读/写)
\label{fielddesc:I2CCOMMAND10DONE}\item [I2C\_COMMAND10\_DONE] 在 I2C 主机模式下完成命令 10 时，该位翻转为高电平。 (读/写)
\end{reglist}\end{regdesc}
\end{register}


\begin{register}{H}{I2C\_COMD11\_REG}{\hexprefix{}0084}\label{regdesc:I2CCOMD11REG}
\regfield{I2C\_COMMAND11\_DONE}{1}{31}{{0}}%
\regfield{(reserved)}{17}{14}{00000000000000000}%
\regfield{I2C\_COMMAND11}{14}{0}{{\hexprefix{}00}}%
\reglabel{Reset}\regnewline%
\begin{regdesc}\begin{reglist}
\label{fielddesc:I2CCOMMAND11}\item [I2C\_COMMAND11] 命令 11 的内容。 该命令包括三个部分：
op\_code 为命令，0：RSTART；1：WRITE；2：READ；3：STOP；4：END。
byte\_num 表示需发送或接收的字节数。
ack\_check\_en、ack\_exp 和 ack 用于控制 ACK 位。 参阅 I2C cmd 结构获取更多
信息。 (读/写)
\label{fielddesc:I2CCOMMAND11DONE}\item [I2C\_COMMAND11\_DONE] 在 I2C 主机模式下完成命令 11 时，该位翻转为高电平。 (读/写)
\end{reglist}\end{regdesc}
\end{register}


\begin{register}{H}{I2C\_COMD12\_REG}{\hexprefix{}0088}\label{regdesc:I2CCOMD12REG}
\regfield{I2C\_COMMAND12\_DONE}{1}{31}{{0}}%
\regfield{(reserved)}{17}{14}{00000000000000000}%
\regfield{I2C\_COMMAND12}{14}{0}{{\hexprefix{}00}}%
\reglabel{Reset}\regnewline%
\begin{regdesc}\begin{reglist}
\label{fielddesc:I2CCOMMAND12}\item [I2C\_COMMAND12] 命令 12 的内容。 该命令包括三个部分：
op\_code 为命令，0：RSTART；1：WRITE；2：READ；3：STOP；4：END。
byte\_num 表示需发送或接收的字节数。
ack\_check\_en、ack\_exp 和 ack 用于控制 ACK 位。 参阅 I2C cmd 结构获取更多
信息。 (读/写)
\label{fielddesc:I2CCOMMAND12DONE}\item [I2C\_COMMAND12\_DONE] 在 I2C 主机模式下完成命令 12 时，该位翻转为高电平。 (读/写)
\end{reglist}\end{regdesc}
\end{register}


\begin{register}{H}{I2C\_COMD13\_REG}{\hexprefix{}008C}\label{regdesc:I2CCOMD13REG}
\regfield{I2C\_COMMAND13\_DONE}{1}{31}{{0}}%
\regfield{(reserved)}{17}{14}{00000000000000000}%
\regfield{I2C\_COMMAND13}{14}{0}{{\hexprefix{}00}}%
\reglabel{Reset}\regnewline%
\begin{regdesc}\begin{reglist}
\label{fielddesc:I2CCOMMAND13}\item [I2C\_COMMAND13] 命令 13 的内容。 该命令包括三个部分：
op\_code 为命令，0：RSTART；1：WRITE；2：READ；3：STOP；4：END。
byte\_num 表示需发送或接收的字节数。
ack\_check\_en、ack\_exp 和 ack 用于控制 ACK 位。 参阅 I2C cmd 结构获取更多
信息。 (读/写)
\label{fielddesc:I2CCOMMAND13DONE}\item [I2C\_COMMAND13\_DONE] 在 I2C 主机模式下完成命令 13 时，该位翻转为高电平。 (读/写)
\end{reglist}\end{regdesc}
\end{register}


\begin{register}{H}{I2C\_COMD14\_REG}{\hexprefix{}0090}\label{regdesc:I2CCOMD14REG}
\regfield{I2C\_COMMAND14\_DONE}{1}{31}{{0}}%
\regfield{(reserved)}{17}{14}{00000000000000000}%
\regfield{I2C\_COMMAND14}{14}{0}{{\hexprefix{}00}}%
\reglabel{Reset}\regnewline%
\begin{regdesc}\begin{reglist}
\label{fielddesc:I2CCOMMAND14}\item [I2C\_COMMAND14] 命令 14 的内容。 该命令包括三个部分：
op\_code 为命令，0：RSTART；1：WRITE；2：READ；3：STOP；4：END。
byte\_num 表示需发送或接收的字节数。
ack\_check\_en、ack\_exp 和 ack 用于控制 ACK 位。 参阅 I2C cmd 结构获取更多
信息。 (读/写)
\label{fielddesc:I2CCOMMAND14DONE}\item [I2C\_COMMAND14\_DONE] 在 I2C 主机模式下完成命令 14 时，该位翻转为高电平。 (读/写)
\end{reglist}\end{regdesc}
\end{register}


\begin{register}{H}{I2C\_COMD15\_REG}{\hexprefix{}0094}\label{regdesc:I2CCOMD15REG}
\regfield{I2C\_COMMAND15\_DONE}{1}{31}{{0}}%
\regfield{(reserved)}{17}{14}{00000000000000000}%
\regfield{I2C\_COMMAND15}{14}{0}{{\hexprefix{}00}}%
\reglabel{Reset}\regnewline%
\begin{regdesc}\begin{reglist}
\label{fielddesc:I2CCOMMAND15}\item [I2C\_COMMAND15] 命令 15 的内容。 该命令包括三个部分：
op\_code 为命令，0：RSTART；1：WRITE；2：READ；3：STOP；4：END。
byte\_num 表示需发送或接收的字节数。
ack\_check\_en、ack\_exp 和 ack 用于控制 ACK 位。 参阅 I2C cmd 结构获取更多
信息。 (读/写)
\label{fielddesc:I2CCOMMAND15DONE}\item [I2C\_COMMAND15\_DONE] 在 I2C 主机模式下完成命令 15 时，该位翻转为高电平。 (读/写)
\end{reglist}\end{regdesc}
\end{register}


\begin{register}{H}{I2C\_DATE\_REG}{\hexprefix{}00F8}\label{regdesc:I2CDATEREG}
\regfield{I2C\_DATE}{32}{0}{{\hexprefix{}19052000}}%
\reglabel{Reset}\regnewline%
\begin{regdesc}\begin{reglist}
\label{fielddesc:I2CDATE}\item [I2C\_DATE] 版本控制寄存器。 (读/写)
\end{reglist}\end{regdesc}
\end{register}


